% problem-000439 1 金属盐热分解转化
\section*{1. 金属盐热分解转化}
M 为常⻅⾦属，不溶于盐酸和稀硫酸，质地坚韧，富延展性。将 M 溶于硝酸后蒸发、浓缩、冷却得到 A的⽔合盐。A受热分解得到⿊⾊固体 $\mathbf { B } _ { \circ }$ 。B溶于盐酸，经蒸发、浓缩、冷却得到绿⾊晶体 $\mathbf { C } _ { \circ }$ 。B溶于稀硫酸后经过蒸发、浓缩、冷却得到蓝⾊晶体D。D在 $2 7 0 ~ ^ { \circ } \mathrm { C }$ 恒温下⽣成⽩⾊粉末 $\mathbf { E } ,$ ，E经 $6 0 0 ~ ^ { \circ } \mathrm { C }$ 恒温后⽣成 $\mathbf { B } _ { \circ }$

\subsection*{1-1}
写出A、B、C和E的化学式。

\subsection*{1-2}
将E溶于⽔，加⼊过量的碘化钾溶液充分反应，其后缓慢滴加硫代硫酸钠溶液⾄过量。简述反应过程中发⽣的现象，写出相关反应的离⼦⽅程式。

\subsection*{1-3}
A的熔点较低，真空时易升华，这与⼀般离⼦晶体的性质不相符。简述理由。

\subsection*{1-4}
M 在⼀定温度下与氧作⽤得到固体 $\mathbf { G } _ { \circ }$ 。G 属⽴⽅晶系，其正当晶胞如下图所⽰，密度为 $6 . 0 0 \ \mathrm { g \ c m ^ { - 1 } c }$ 。通过计算给出M‒O的距离。

（图：）
